\begin{table}[t]
\centering
\small
\setlength{\tabcolsep}{3.5pt}
\begin{tabular}{clccc}
\toprule
ID & 체크포인트 & PT & IT & Traj \\
\midrule
C0 & Base & -- & -- & -- \\
C1 & Short PT & Short & -- & -- \\
C2 & Long PT & Long & -- & -- \\
C3 & Base + IT & -- & \checkmark & -- \\
C4 & Short PT + IT & Short & \checkmark & -- \\
C5 & Long PT + IT & Long & \checkmark & -- \\
C6 & Base + IT + Traj & -- & \checkmark & \checkmark \\
C7 & Short + IT + Traj & Short & \checkmark & \checkmark \\
C8 & Long + IT + Traj & Long & \checkmark & \checkmark \\
\bottomrule
\end{tabular}
\caption{\textbf{평가할 아홉 개 체크포인트.} 중간 체크포인트를
보존하여 PT 효과가 나타나는 지점과 후속 학습 이후의 유지 여부를
분리한다.}
\label{tab:checkpoints}
\end{table}
